\ficha{Franco (1988), Assimetrias sistêmicas sob o padrão ouro}{franco1988}

\begin{identificacao}
FRANCO, Gustavo H. B. \textit{Assimetrias sistêmicas sob o padrão ouro}. Rio de
Janeiro: Departamento de Economia, Pontifícia Universidade Católica do Rio de
Janeiro, 1988. (Texto para Discussão, n. 184). 22 p.

Ensaio de resenha interpretativa, em português, lido na íntegra: 15 páginas de
texto corrido, quatro páginas de notas e três de referências, mais um gráfico e
três tabelas em folhas prévias sem paginação. Versão lida: cópia digitalizada
depositada no EconStor (handle 10419/186431), cujas duas primeiras páginas são
capa do repositório. O trabalho integra a pesquisa ``O Brasil e a Economia
Internacional, 1870-1914'' e anuncia um segundo ensaio, sobre o caso
brasileiro, que não está neste texto.
\end{identificacao}

\bloco{Problema e tese}

A maior instabilidade macroeconômica dos países periféricos entre 1870 e 1914 é
efeito de características próprios desses países, uma \textit{self inflicted
wound}, ou é produto de assimetrias embutidas nos mecanismos de ajustamento da
economia internacional? Franco percorre o segundo caminho e sustenta que a
estabilidade do centro e a instabilidade da periferia são as duas faces do
mesmo arranjo.

\chave{o argumento interessante a explorar seria o de que uma contrapartida à
estabilidade experimentada no centro seria a instabilidade na periferia, ou que
o primeiro podia exportar instabilidade à segunda}{p. 1}

\bloco{Argumento}

\begin{enumerate}[leftmargin=*,nosep]
\item A questão da ``gerência'' do sistema é deslocada da intenção para o
efeito. Não havia missão declarada do Banco da Inglaterra de administrar a
economia internacional, e Schumpeter é citado sobre a discrição deliberada de
seus dirigentes (p. 3). O que importa é a funcionalidade sistêmica da conduta
rotineira do Banco, e o fato de que a estratégia britânica de reconstrução do
padrão ouro após 1918 supunha que ele tinha sido o gestor do sistema antes de
1914 (p. 3).

\item A distribuição geográfica da liquidez é assimétrica. As reservas
metálicas do Banco da Inglaterra eram pequenas, inferiores às de vários países
europeus, e Japão e Índia mantinham reservas em libras maiores que as reservas
metálicas britânicas. O sistema era de fato um ``padrão libra'', e não tinha
como fundamento básico a conversibilidade (p. 5). A vulnerabilidade era mais
aparente que real porque parte substancial dos saldos em libras era inerte, a
Índia tinha autonomia reduzida sobre suas reservas, e a taxa de desconto
funcionava como \textit{deus ex machina} em todas as crises do período, embora
fosse apenas um ajuste na margem (p. 6).

\item A transmissão de crises é hierárquica. Londres emitia sinais que
pressionavam primeiro os centros mais próximos e daí irradiavam em ondas para
as extremidades (p. 7). A leitura que Franco prefere é a de transferência
temporária de ativos líquidos da periferia para o centro, ou seja, os centros
privilegiados transferiam o ônus do ajuste a um aumento da preferência pela
liquidez para as extremidades mais frágeis do sistema (p. 8).

\item O balanço de pagamentos britânico consolida a assimetria. A conta
corrente era cronicamente superavitária apesar do déficit comercial, graças a
fretes e rendas de investimento, e a contrapartida foi uma exportação de
capital que atingiu em média 5,2\% do PNB ao ano (p. 11). Credor e devedor não
têm o mesmo repertório: desacelerar exportações de capital é mais fácil que
obter novos empréstimos ou ajustar a conta corrente.

\chave{há uma assimetria óbvia nesse sistema na medida que credores e devedores
possuem capacidades inteiramente diversas para enfrentar desequilíbrios
externos}{p. 11-12}

\item Na versão cíclica, a assimetria se agrava para o exportador primário.
Aceita a correlação negativa entre exportações de capital britânicas e termos
de troca britânicos, o capital chega ao país receptor exatamente quando seus
termos de troca estão favoráveis, gerando ``uma dupla bonança'' (p. 14). Em
seguida a atividade e as importações crescem, os termos de troca pioram, forma-se
o déficit em conta corrente, e a desaceleração das entradas de capital obriga
ao ajuste no pior momento possível.

\chave{o déficit em conta corrente tem de ser eliminado via ajustes na renda ou
na taxa de câmbio justamente quando os termos de troca estão em seu pior
momento}{p. 14}

\item O investimento doméstico, que na Grã-Bretanha amortece o ciclo por estar
negativamente correlacionado com as exportações de capital (p. 13), não pode
cumprir esse papel na periferia, porque a necessidade de eliminar o déficit não
o permite (p. 14). Triffin, citado ao fim, acrescenta a dependência dos
devedores em um ou poucos produtos primários de exportação, sujeitos a amplas
variações de preços e quantidades (p. 15).
\end{enumerate}

\bloco{Conceitos e definições operacionais}

\textbf{Assimetria sistêmica}: propriedade do mecanismo de ajustamento, não do
país periférico. Operacionalmente, a capacidade desigual de credores e
devedores de adiar o ajuste (p. 11-12).

\textbf{Padrão libra}: regime internacional cuja base não é a conversibilidade
da libra em ouro, nem o estoque metálico do Banco da Inglaterra, mas a posição
central da Grã-Bretanha no comércio e nas finanças (p. 5 e p. 15).

\textbf{Economia líder}: aquela cuja moeda é moeda internacional de reserva, o
que lhe permite, na expressão de Lindert que Franco adota, o ``déficit sem
lágrimas'' (p. 12).

\textbf{Transferência do ônus}: hipótese de que a resposta do centro a um choque
de preferência pela liquidez desloca o ajuste para as extremidades, exportando
crises financeiras à periferia (p. 8).

\textbf{Áreas de influência} ou ``suseranias'', na expressão de De Cecco:
ligações privilegiadas que fazem o capital percorrer canais forjados pela
experiência e pela informação acumulada (p. 6).

Franco não define padrão-ouro como regra nem trabalha com credibilidade ou
compromisso. Esse vocabulário está ausente.

\bloco{Evidência e método}

Não há dado novo nem estimação própria. O ensaio organiza evidência de
terceiros: a tabela 1 traz a composição dos estoques de reservas internacionais
em 1913, em milhões de dólares, adaptada de Lindert (1969); a tabela 2, a
importação líquida de ouro pela Grã-Bretanha em 1903-1910, de Ford (1964, p.
34); a tabela 3, o balanço de pagamentos britânico em 1907, 1910 e 1913, de
Sayers (1976); o gráfico 1 traz variáveis macroeconômicas britânicas de
1870-1913 em desvios sobre médias móveis de nove anos, sem crédito de fonte na
legenda. A contestação de Ford ao argumento da transferência do ônus é
respondida por releitura da mesma tabela 2, e Franco reconhece que a disputa só
se resolveria com estatísticas completas de movimentos de ouro que inexistem, e
que as disponíveis são de péssima qualidade (p. 9). A ``relação estatística
entre exportações e exportações de capital apresentada abaixo'' anunciada na
página 12 não aparece no texto, e remete em nota a Ford (1962, p. 73).

\bloco{Agentes e vertentes}

Franco se filia à leitura estruturalista da assimetria, com Triffin, Lindert e
De Cecco, contra a visão clássica do automatismo, cuja linhagem ele traça de
Hume e Ricardo a Mill, Taussig e os ``neoclássicos de Harvard'' (John H.
Williams sobre a Argentina, Graham sobre os EUA, Viner sobre o Canadá, White
sobre a França). Keynes fornece a imagem do Banco da Inglaterra como maestro de
uma orquestra internacional, e Schumpeter a observação sobre a discrição do
Banco. Bloomfield e Nurkse sustentam que os bancos centrais praticavam
neutralização em vez das regras do jogo. Kindleberger fornece o relato de
transmissão de crises em 1888-1893, que inclui o Brasil. Alec Ford é o
contraditor nomeado, com o pânico de 1907. O Cunliffe Report entra como exemplo
de visão ingênua. O Brasil aparece três vezes, e sempre como dado: na tabela 1,
com 89,6 milhões de dólares de reservas, cem por cento em ouro, sob a Caixa de
Amortização; na lista de crises de Kindleberger; e na rubrica ``América do
Sul'' da tabela 2.

\bloco{Críticas e limites}

O texto é programático e o próprio autor o admite. A hipótese da transferência
do ônus tem ``pouca evidência sólida'' (p. 8), e o cenário cíclico da periferia
é apresentado como ``a hipótese mais perversa possível'', com o reconhecimento
de que as histórias individuais podem mostrar padrões bastante divergentes (p.
14). O ensaio não testa nada: rebate Ford sem produzir evidência nova.

Não há tratamento de dívida soberana, de banqueiros, de condicionalidade nem de
política cambial de Estado periférico. Não há teoria de por que uma periferia
adota conversibilidade, e caixa de conversão não é objeto. A perspectiva é a do
balanço de pagamentos britânico, e a periferia entra como contrapartida
agregada, sem agência.

O Brasil não é analisado. A inserção da economia brasileira é explicitamente
adiada ao próximo ensaio (p. 1). Nada aqui sustenta afirmação sobre a Caixa de
Conversão, e o argumento transfere como hipótese a testar, não como medida.

O ensaio é anterior à literatura de credibilidade dos anos 1990, com a qual não
dialoga. Quem tomar Franco como contraponto a Bordo e Kydland faz uma
reconstrução legítima, mas retrospectiva.

\begin{dialogo}
\item Procurar, em 1906-1913, se os jornais atribuem o movimento da taxa a
fatores externos ou a responsabilidade doméstica. Vocabulário: ``taxa de
desconto do Banco da Inglaterra'', ``praça de Londres'', ``mercado monetário
europeu'', ``crise de Nova York''. Testa se o debate brasileiro percebia a
assimetria de Franco, ou se enquadrava a instabilidade como problema de
política interna, o que é a alternativa da \textit{self inflicted wound}.

\item Rastrear a sequência da ``dupla bonança'': se em 1906-1908 os artigos
associam entrada de empréstimos e alta cambial ao bom preço do café, e se em
1913-1914 descrevem a reversão como simultânea, capital escasso e preço em
baixa. Testa se a defesa da Caixa era apresentada como absorção de bonança ou
como proteção contra a reversão, distinção que separa o argumento ortodoxo do
expansionista.

\item Ler a cobertura do pânico de 1907 nos quatro jornais e verificar se o
movimento da Caixa é lido como isolante ou como correia de transmissão do
choque externo. Testa a hipótese de transferência do ônus com material que
Franco diz não existir para a periferia.

\item No debate da custódia do ouro em 1911 e no da agência de Londres,
procurar se algum agente argumenta que reserva depositada em Londres serve de
liquidez à praça inglesa. Vocabulário: ``ouro depositado em Londres'',
``agência de Londres'', ``garantia'', ``depósito''. Testa se a assimetria da
tabela 1, reserva periférica como funding do centro, era percebida na época.
\end{dialogo}

\usonocapitulo{Parte I, como formulação do argumento de assimetria
centro-periferia do ajuste, que o capítulo 2 já invoca. Serve a duas afirmações
precisas: a assimetria está na capacidade desigual de credores e devedores de
adiar o ajuste (p. 11-12), e, na versão cíclica, na coincidência entre seca de
capital e piora dos termos de troca para um exportador primário (p. 14). Serve
também de contraponto à leitura do padrão-ouro como regra de credibilidade,
porque para Franco o sistema anterior a 1914 assentava menos na
conversibilidade que na posição central da Grã-Bretanha (p. 5). Como o texto
não trata do Brasil e adia o caso ao ensaio seguinte, entra como moldura
analítica, nunca como evidência sobre a Caixa de Conversão.}
